\begin{table}[t]
\centering
\small
\caption{Scope of claims and reviewer-facing boundary conditions.}
\label{tab:scope}
\begin{threeparttable}

\begin{adjustbox}{width=\textwidth}
\begin{tabular}{p{0.31\textwidth}p{0.18\textwidth}p{0.43\textwidth}}
\toprule
Question & Answer & Evidence and boundary \\
\midrule
Are confidence/option scores strong for MCQ selective prediction? & Yes & Sub-2B study, seven-model 7B/9B-class audit, fixed-coverage risk, AURC/E-AURC. \\
Do hidden-state probes consistently improve over cheap signals? & Only marginally & Hidden-augmented probes improve AUROC in some cases, but mean gains are small: +0.004 over all 28 cases and +0.008 in the competent subset. \\
Are hidden states useless? & No & PCA hidden-only and layer-localization analyses show secondary signal, especially in late/final layers. \\
Does the claim generalize to open-ended hallucination, RAG, or long-form generation? & No & This is an answerable MCQ forced-choice diagnostic study, not a universal hallucination detector. \\
Is this a scaling law over all SLMs? & No & The scaling audit covers seven accessible 7B/9B-class families, not all 7B--70B models or base/instruction variants. \\
\bottomrule
\end{tabular}
\end{adjustbox}

\end{threeparttable}
\end{table}
